\chapter{Projeto de Pesquisa ou Plano de Trabalho}

O projeto é o documento de planejamento do TCC. Ele organiza as decisões antes da execução e permite que orientador, estudante e curso avaliem viabilidade, coerência e riscos.

Sua apresentação deve considerar a ABNT NBR 15287:2025 e as exigências do curso. O projeto não é uma versão abreviada do TCC final: ele descreve uma pesquisa planejada, enquanto o trabalho final registra e analisa o percurso efetivamente realizado.

\section{Função do projeto}

O projeto serve para responder antecipadamente a quatro perguntas: o que será estudado, por que será estudado, como será estudado e em quanto tempo será realizado. Ele não é uma formalidade isolada; é a base do trabalho final.

Mesmo quando o curso não exige entrega formal de projeto, o estudante deve elaborar um plano mínimo. Trabalhos iniciados sem planejamento tendem a enfrentar problemas de recorte, método e prazo.

O projeto não é uma versão abreviada do TCC final. Ele registra o que se pretende fazer e pode ser ajustado de forma justificada durante a execução. O trabalho final, por sua vez, relata o percurso efetivamente realizado, analisa os resultados e apresenta a parte conclusiva.

\section{Organização do projeto}

Conforme a ABNT NBR 15287:2025, o projeto é organizado em elementos pré-textuais, textuais e pós-textuais. O conteúdo e a extensão variam conforme a natureza da pesquisa e as regras do curso.

\subsection{Elementos pré-textuais}

A folha de rosto é obrigatória no projeto. Capa, listas e sumário são usados conforme a necessidade e a orientação institucional. O modelo do curso deve indicar quais informações de identificação e vínculo acadêmico serão exigidas.

\subsection{Elementos textuais}

Um projeto de TCC pode conter:

\begin{itemize}
  \item título provisório;
  \item tema e delimitação;
  \item problema de pesquisa;
  \item hipótese ou pressuposto, quando couber;
  \item objetivo geral e objetivos específicos;
  \item justificativa;
  \item referencial teórico inicial;
  \item metodologia;
  \item cronograma;
  \item cronograma.
\end{itemize}

As seções não precisam receber exatamente esses títulos, mas o texto deve permitir reconhecer tema, problema, objetivos, justificativa, base teórica, método, recursos quando necessários e planejamento temporal.

\subsection{Elementos pós-textuais}

As referências são obrigatórias. Glossário, apêndices, anexos e índice são incluídos quando necessários. Instrumentos de coleta elaborados pelo estudante, como roteiros e questionários, são apêndices; documentos de terceiros são anexos.

\section{Referencial teórico inicial}

O referencial teórico inicial apresenta autores, conceitos e debates que sustentam a pesquisa. Não precisa esgotar a literatura, mas deve demonstrar que o estudante conhece o campo mínimo necessário para iniciar o trabalho.

É recomendável evitar duas fragilidades: listar autores sem articulação ou transformar o referencial em resumo de obras. O referencial deve construir base conceitual para o problema.

\section{Metodologia no projeto}

A metodologia deve explicar o tipo de pesquisa, os materiais ou participantes, os procedimentos, os critérios de análise e os cuidados éticos. Quanto mais claro for o método no projeto, menor a chance de retrabalho durante a execução.

Quando houver coleta com pessoas, dados sensíveis, intervenção, prontuários, imagens, entrevistas ou observação de campo, a necessidade de avaliação ética deve ser verificada com antecedência.

\section{Cronograma e riscos}

O cronograma deve considerar etapas reais e riscos previsíveis. Se o trabalho depende de autorização externa, resposta de participantes, acesso a documentos ou aprovação ética, essas etapas precisam aparecer no planejamento.

\caixaatencao{Não deixe revisão bibliográfica, normalização e conferência de referências para a última semana. Essas etapas costumam revelar inconsistências que exigem reescrita.}

\section{Recursos e viabilidade}

Quando necessários, devem ser descritos recursos humanos, materiais, tecnológicos e financeiros, bem como formas de acesso a campo, bases, equipamentos ou documentos. Dependências críticas precisam aparecer no planejamento e no cronograma.
